\section{A low-power adjoint bound for unitary natural dimension}
\label{sec:unitary-adjoint-rank}

\paragraph{Accepted status and shared origin.}
D-TD69, question Q2fy, is an accepted effective natural-dimension bound.
The complete written mathematics, exact checker, and first successful
output pass author, independent whole-packet and root review.
The immutable author freeze and all earlier working histories are retained.
The parent proposed the cyclic low-support subgroup and an induced
adjoint module, and the author supplied the inner-diagonal trace
estimate and the two-index presentation replacement. The finite
ordinary-Schreier relator argument is the shared, accepted D-TD64
argument and the separate Sixfolds relator lemma. No independent
origin is claimed for those inputs. The full unitary outer-centralizer
argument is the accepted D-TD68 argument; the present proof repeats
the exact quantitative consequence with a weaker degree floor.
It uses neither the D-TD66 genus comparison nor a level restriction.

\paragraph{Constants and one computable finite-presentation function.}
Let \(M\ge80\) be an integer and \(q=p^f\ge2\) a prime power.
Put
\[
 L=\lceil\log_2(32M)\rceil,\qquad T=10^{33}L,
 \qquad \epsilon=(16M)^{-1},\qquad \sigma=10^{-29},
 \qquad
 O= p^{\lceil\log_p(T+1)\rceil}
       \operatorname{lcm}_{1\le h\le T+1}(q^{2h}-1).
 \tag{UAR.1}
\]
Define \(D(M)\ge1\) as follows. Enumerate every integer
\(1\le r\le3M\) and every tuple \((b_1,\ldots,b_r)\) with
\(1\le b_i\le M\) and \(\sum_i(1-1/b_i)<2\).
For its presentation
\[
 \Gamma_b=\langle x_1,\ldots,x_r:
                x_1\cdots x_r=1,\ x_i^{b_i}=1\rangle,
 \tag{UAR.2}
\]
use the accepted D-TD64 spherical-presentation and ordinary-Schreier
algorithm. Eliminate \(x_r\), choose the actual breadth-first Cayley
transversal of \(\Gamma_b\), and find free-word certificates expressing
every ordinary Schreier generator of the kernel as a product of
conjugates of the power relators and their inverses. Let its total
certificate area be the sum of these word areas. Fix a deterministic
length-then-lexicographic certificate enumeration throughout. The maximum of one
and all these finitely many totals is \(D(M)\).

This is a terminating integer algorithm. After deleting entries one,
strict sphericity leaves at most three orders. The two-order groups are
cyclic, and the three-order signatures are \((2,2,c)\),
\((2,3,3)\), \((2,3,4)\), \((2,3,5)\). Hence
\(|\Gamma_b|\le\max(2M,60)=2M\). The accepted covering-space proof
identifies the entire presentation kernel, not only a finite quotient.
Alternatively the relator algorithm finds the actual maximal marked
group table of order at most \(2M\), then searches for the kernel
certificates; the finite-presentation theorem guarantees termination.
No closed-form area bound or execution of this enormous algorithm is
claimed.

Set
\[
 R(q,M)=\max\{4T,\ 2M+1,
                 2OT(1+3M^2)D(M)+2\}.
 \tag{UAR.3}
\]

\begin{theorem}[D-TD69]
Let \(C\to\mathbf P^1\) be an actual connected finite \(J\)-Galois
cover. Let \(V\) be a faithful irreducible complex coefficient module,
\(e=\dim V\), with its genuine generating product-one inertia tuple
\(T_1,\ldots,T_r\). Suppose
\[
 r\ge12,\qquad \sum_i\operatorname{rank}(T_i-I)=2e+6,
 \qquad \operatorname{rank}(T_i-I)\ge e/M\quad(1\le i\le r).
 \tag{UAR.4}
\]
The last assumption follows, in particular, from a whole-source
projective support floor \(1/M\): every genuine scalar inertia costs
its entire degree.

Suppose \(V=W\otimes U\) is preserved factorwise and its actual finite
projective image \(P\) on \(W\) contains and normalizes the projective
image of a faithfully represented irreducible central product \(E_1\cdots E_t\).
The factors are commuting quasisimple groups with central
intersections, and \(P\) acts transitively on them. Suppose each
\(E_j\) is an actual central quotient of the named group
\(\operatorname{SU}_d(q)\), with \(d\ge120\).
There is no assumption on the entire external partner \(U\), no
bound on \(t\), and no character-level restriction. Then
\[
                              d<R(q,M).
 \tag{UAR.5}
\]
An exceptional central cover is not included merely from the name of
its simple quotient. No bound on original scalar inertia orders is
asserted.
\end{theorem}

\subsection*{1. The exact represented normalizer and a degree floor}

Write \(S=\operatorname{PSU}_d(q)\),
\(A=\operatorname{Aut}(S)\), \(B=\operatorname{PGU}_d(q)\), and
\(c=[A:B]\le2f\). The named SU group is quasisimple in this range.
The component modules \(W_j\) have a common degree \(a>1\).
Perfectness kills scalar discrepancies in projective normalization.
If a represented normalizer induces the identity on every simple
component, it centralizes the actual perfect core, so Schur's lemma
makes it scalar. Thus the actual group, not a substituted larger
normalizer, embeds as
\[
             K=S^t\ \triangleleft\ P\ \le A\wr S_t.
 \tag{UAR.6}
\]
Steinberg's complete twisted-type automorphism proof
\cite[pp.~607--608 and complete Section~7, pp.~613--614]{PFL:Steinberg1960}
supplies
\(B\triangleleft A\) and a semilinear model
\(X\mapsto H X^\tau H^{-1}\), with \(H\in\operatorname{GU}_d(q)\)
and \(\tau\in\operatorname{Gal}(\mathbf F_{q^2}/\mathbf F_p)\).
The graph operation is the field map \(X\mapsto X^q=X^{-T}\).
These are precisely the named-cover/perfectness interfaces written
in D-TD68; no incompatible automorphism is asserted to stabilize a
given character.

The accepted minimum-degree input CFD.7 gives every nontrivial
ordinary irreducible \(S\)-module degree at least \(q^{d/6}\).
The ordinary module \(\operatorname{End}_{\mathbf C}(W_j)\) factors
through \(S\), has exactly one invariant line, and is not trivial.
Consequently
\[
                      a^2-1\ge q^{d/6},\qquad a>q^{d/12}.
 \tag{UAR.7}
\]
This argument does not confuse the minimum degree of the simple
quotient with that of an arbitrary central-cover representation.
The latter degree floor is all that will be needed.

We may assume \(d\ge R(q,M)\) and derive a contradiction. In
particular \(d\ge4T\), so \(e\ge a^t>2^{d/12}>6M\).
Equation (UAR.4) therefore gives \(r<3M\).

\subsection*{2. Trace decay for every bad component operator}

The accepted LT bound for an inflated SU character
\cite[Theorem~5.5 and complete proof, author pp.~25--26]{UNW:LT2024} is
\( |\chi(x)|/a\le a^{-\sigma\operatorname{supp}(x)/d}\).
Together with (UAR.7), this is at most
\(q^{-\sigma\operatorname{supp}(x)/12}\).
The support here is geometric natural support; central lift choices
do not affect its value or the absolute character ratio.

First consider an actual compatible inner-diagonal local operator
inducing conjugation by \(H\in\operatorname{GU}_d(q)\), with
\(s=\operatorname{supp}(H)\ge64\). The complete statement and proof
of LT Proposition 6.1, author pp.~27--28, apply to every such \(H\)
\cite[Proposition~6.1 and its complete proof]{UNW:LT2024}
and give \(|H^{\operatorname{SU}}|\ge|\operatorname{SU}|^{s/(2d)}\).
There are at most \(q+1\) projective scalar-discrepancy fibers, each
a coset of the ordinary SU centralizer of \(H\). Using
\(|\operatorname{SU}_d(q)|\ge(2/3)q^{d^2-1}\), we obtain
\[
 \frac{|C_S(\alpha_H)|}{|S|}
 \le(q+1)|\operatorname{SU}|^{-s/(2d)}
 \le4q^{1-ds/2+s/(2d)}.
 \tag{UAR.8}
\]
No extension of \(\chi\) from SU to GU is required for this bound.

Put \(R_s=\lfloor s/32\rfloor\), so \(R_s\ge s/64\).
At most \(q^{4dR_s+2}\) simple-group elements have support below
\(R_s\): their dominant eigenvalue is rational, and a scalar-shifted
natural lift is overcounted by a product of a \(d\)-by-\(R_s\) and an
\(R_s\)-by-\(d\) matrix over \(\mathbf F_{q^2}\).
Every nonempty commutator-map fiber has size \(|C_S(\alpha_H)|\).
Thus its low-support fraction is at most
\(q^{-3ds/8+11/2}\le q^{-ds/4}\), since \(ds\ge44\).
Schur averaging, with a common finite unitary lift of the actual
projective normalizer, now gives
\[
 \left|\frac{\operatorname{tr}A_H}{a}\right|^2
       \le q^{-ds/4}+q^{-\sigma s/768}
       \le2q^{-\sigma s/768}.
 \tag{UAR.9}
\]

For a local automorphism outside \(B\), the complete centralizer
argument of D-TD68 gives projective fixed fraction
\(3q^{-d^2/3+2}\). For clarity, a field part of order at least three
has fixed matrix-space size at most \(q^{2d^2/3}\) per scalar fiber.
A field part of order two is the inverse-transpose coset; the
Gow--Vinroot centralizer correspondence
\cite[complete Section~2, author pp.~3--6]{UPL:GowVinroot2008}
transfers it to the isometry
group of a nonsingular bilinear form over \(\mathbf F_q\), whose
column count is at most \(q^{d(d+1)/2}\). These statements hold in
both characteristics and for arbitrary coset-element orders.
The extra projective fibers contribute at most \(q+1\).
At \(R_d=\lfloor d/32\rfloor\), the same low-support count and
Schur average, now using the weaker (UAR.7), give
\[
 \left|\frac{\operatorname{tr}A_\alpha}{a}\right|^2
       \le q^{-d^2/6}+q^{-\sigma d/768}
       \le2q^{-\sigma d/768}\qquad(\alpha\notin B).
 \tag{UAR.10}
\]
Only the full-normalizer proof of D-TD68 is invoked here, not its
level-two or genus-comparison premise. Its supplied proofs use
Steinberg, Gow--Vinroot and the elementary unitary order calculation
\cite[pp.~114--118]{OS:Taylor1992}.

Because \(\sigma T/768=(10^4/768)L>13L\), (UAR.9)--(UAR.10)
are at most \(\epsilon^2\) when respectively \(s>T\) or
\(d\ge4T\). A component-moving tensor operator has normalized
absolute trace at most \(1/a<\epsilon\): its exact cycle-trace
formula has a factor \(a^{-I}\), with \(I\ge1\), while all unitary
cycle traces are at most their degrees. Consequently every actual
projective element outside the following set has whole-layer
normalized absolute trace at most \(\epsilon\):
it fixes every component, acts inner-diagonally on each, and has
natural support at most \(T\) on each.
Arbitrary scalar phases and the entire external partner preserve
these absolute-trace upper bounds.

\subsection*{3. Cyclic low-support subgroups and full original periods}

Let \(g_i\) be the actual image of the original \(T_i\) in \(P\).
Inside the finite cyclic group \(\langle g_i\rangle\), let \(E_i\)
be the subgroup generated by all elements of the set just described.
This definition includes the identity, requires every component to
be fixed, and uses the actual natural action on every component.
Put \(d_i=[\langle g_i\rangle:E_i]\).

Every generating low-support element has projective order dividing
\(O\). Indeed, when \(d>2T\), its dominant eigenvalue on each
natural component is rational. After scalar normalization, the
matrix \(H\) has \(\operatorname{rank}(H-I)\le T\). Its image
\((H-I)\mathbf F_{q^2}^d\) is \(H\)-invariant; the minimal
polynomial of the restriction has degree at most \(T\), so
\((X-1)\) times that polynomial annihilates \(H\).
Every semisimple eigenvalue therefore has extension degree at most
\(T+1\), and every unipotent block has size at most \(T+1\).
The expression \(O\) in (UAR.1) is a common multiple of their orders.
It also bounds the order of a componentwise tuple, independently
of \(t\). A subgroup of a cyclic group generated by elements with
orders dividing \(O\) has order dividing \(O\). Hence
\[
                 |E_i|\mid O,\qquad E_i=\langle g_i^{d_i}\rangle.
 \tag{UAR.11}
\]

Retain the full order \(L_i\) of the original coefficient operator
\(T_i\). Its projective layer order divides \(L_i\). The original
tensor factors can be normalized by reciprocal unit scalars so
that their \(L_i\)-th powers are identity. Exactly a fraction
\(1/d_i\) of these full-period powers lie in \(E_i\); all powers
outside \(E_i\) have layer absolute trace at most \(\epsilon\).
Fourier averaging, at every eigenvalue target and with the entire
partner trace bounded by its degree, gives
\[
 \frac{\operatorname{rank}(T_i-I)}e
       \ge (1-1/d_i)(1-\epsilon).
 \tag{UAR.12}
\]
No order of an original scalar lift has been replaced by a projected
order, and no source budget is assigned to \(W\) alone.

Let \(h\) count the indices with \(d_i>1\), and
\(S_0=\sum_i(1-1/d_i)\). Five such indices would cost at least
\(5(1-\epsilon)/2>2+1/M>2+6/e\), so \(h\le4\).
At least eight of the original genuine marks remain. Charging their
whole-source cost \(1/M\) separately gives
\[
 (1-\epsilon)S_0+8/M\le2+6/e<2+1/M,
 \qquad S_0<2-6/M.
 \tag{UAR.13}
\]
The last elementary inequality follows from
\(2-7/M<(2-6/M)(1-1/(16M))\).
Four positive indices have \(S_0\ge2\), so in fact \(h\le3\).
The omitted marks are genuine original inertia, including
projectively scalar or layer-trivial entries; they are not erased.

\subsection*{4. An honest induced adjoint module}

Put \(F=\mathbf F_{q^2}\) and \(Z_0=\operatorname{End}_F(F^d)\).
The group \(B\) acts on \(Z_0\) by conjugation, independently of
its scalar natural lifts. Use the honest finite-dimensional
\(F[A]\)-module \(\operatorname{Ind}_B^A Z_0\); this construction
does not presume a semilinear splitting of \(A\) or an absolutely
irreducible adjoint quotient. Form the direct sum of its \(t\)
component copies, with the natural wreath action, and restrict to
the actual \(P\). Call this module \(\mathcal A\), and write
\(\rho:P\to\operatorname{GL}_F(\mathcal A)\). Then
\[
 \dim_F\mathcal A=ctd^2,\qquad
 \dim_F[K,\mathcal A]=ct(d^2-1).
 \tag{UAR.14}
\]
To prove the second identity, restriction to \(B\) is a direct sum
of \(c\) automorphism twists of \(Z_0\). On each, the natural
SU preimage is absolutely irreducible: the determinant-one norm-one
diagonal group has distinct coordinate characters, and even
permutation matrices are transitive. The scalar commutant therefore
has dimension one. The trace pairing on the full matrix algebra is
nondegenerate in every characteristic and identifies the orthogonal
complement of the commutator span with that commutant. The span
thus has dimension \(d^2-1\), even if \(p\mid d\).
Sum over the frames and components. Neither simplicity nor
irreducibility of the adjoint module is used.

For an inner-diagonal natural matrix of support at most \(T\),
its dominant scalar is rational and
\(\operatorname{rank}(\operatorname{Ad}H-I)\le2dT\), by writing
\(H\otimes H^{-T}-I\) as the sum of two scalar-shift terms.
Every conjugate frame preserves geometric support. Each generating
low-support element consequently has \(\rho\)-moving rank at most
\(2ctdT\). A shortest word in this generating set and its inverses
has length at most \(|E_i|-1\le O\). The rank product inequality
therefore gives
\[
         \operatorname{rank}(\rho(g_i^{d_i})-I)
                  \le Z:=2ctdOT.
 \tag{UAR.15}
\]
Inverses have the same support and rank, so no generating-set
symmetry assumption is missing. Explicitly, a shortest Cayley path
has no repeated vertex; if equal steps are grouped into integer
exponents, their total absolute coefficient is no greater than the
path length. No uncontrolled Bezout coefficients are used.

\subsection*{5. Bounded replacement exponents}

If \(h=3\), sort the three nontrivial indices. Strict sphericity
gives one of the classical spherical triples. In the dihedral case
\((2,2,c')\), (UAR.13) gives \(1/c'>6/M\), hence \(c'<M/6\).
The other three triples have all orders at most five. Thus every
\(d_i\le M\); retain all these exponents.

If \(h=2\), write the two positive indices as \(d_1\le d_2\),
relabeling their positions for this argument. Equation (UAR.13)
gives \(1/d_1+1/d_2>6/M\), hence \(\delta=d_1<M/3\).
Retain exponent \(\delta\) on the first entry and replace the
second exponent by \(\delta\). All other exponents are one.
To justify this replacement, cyclically rotate the original
product-one tuple, without changing any local bounds, and write
\(g_1 A g_2 B'=1\), with \(A,B'\) products of the other entries.
Every such entry has \(d_i=1\) and \(\rho\)-moving rank at most
\(Z\). Thus
\[
 \operatorname{rank}\bigl(\rho(g_2)-\rho(g_1)^{-1}\bigr)
       \le(r-2)Z,\qquad
 \operatorname{rank}\bigl(\rho(g_2^{\delta})-I\bigr)
       \le\{1+\delta(r-2)\}Z.
 \tag{UAR.16}
\]
The second inequality is the exact noncommutative telescoping
identity for \(X^{\delta}-Y^{\delta}\), followed by (UAR.15).
It does not assert that \(g_2^{\delta}\) belongs to \(E_2\),
has small natural support, or has bounded order. The resulting
two-order presentation is cyclic of order \(\delta\).

If \(h=1\), replace its sole positive exponent by one. The
product-one relation bounds this entry's \(\rho\)-moving rank by
\((r-1)Z\). If \(h=0\), keep all exponents one.
In all four cases we have exponents \(1\le b_i\le M\), a
spherical presentation (UAR.2), and
\[
 \operatorname{rank}(\rho(g_i^{b_i})-I)
                 \le (1+3M^2)Z\quad(1\le i\le r).
 \tag{UAR.17}
\]
The large original projected orders, if any, have not been bounded
or silently substituted into the finite signature.

\subsection*{6. Ordinary relator generation and the core dimension}

Let \(\mathcal F\) be the free group after eliminating the original
product-one relation, and \(\mathcal R\) the kernel of
\(\mathcal F\to\Gamma_b\). Map \(\mathcal F\) onto the actual
\(P\) using its generating tuple, and let \(N\) be the image of
\(\mathcal R\). Then \(N\triangleleft P\) and
\(|P/N|\le|\Gamma_b|\le2M\).
The ordinary Schreier generators of \(\mathcal R\), not only
normal generators, map to ordinary generators of \(N\).
Each has a certified expression in conjugates of the power relators
and their inverses. Rank is invariant under conjugation and
inversion and subadditive under products. Equations (UAR.17) and
the definition of \(D(M)\) therefore imply
\[
 \dim_F[N,\mathcal A]
             \le D(M)(1+3M^2)Z.
 \tag{UAR.18}
\]
For completeness, the span of the images of the ordinary generators
minus identity equals \([N,\mathcal A]\): their image sum is stable
under each generator, because applying one adds a vector in that
generator's image. It is then stable under inverses and all of
\(N\). This is the exact reason ordinary generation, rather than
normal generation alone, suffices for (UAR.18).

The intersection \(N\cap K\) is a normal subgroup of the direct
product of nonabelian simple groups \(K=S^t\), hence a product of
some components. It is also \(P\)-invariant, and the component
action is transitive. Thus either \(N\cap K=1\) or \(K\le N\).
In the first case \(K\) injects into \(P/N\). The elementary SU
order bound gives
\(|S|\ge q^{d^2-2}/3>d\) for \(d\ge120\), so \(d<2M\),
contrary to (UAR.3).
In the second case, combine (UAR.14), (UAR.15), and (UAR.18) to get
\[
      ct(d^2-1)\le 2ctd\,OT(1+3M^2)D(M),\qquad
      d\le2OT(1+3M^2)D(M)+1.
 \tag{UAR.19}
\]
This again contradicts \(d\ge R(q,M)\), proving the stated
bound. The factors \(c\) and \(t\) cancel. No bad-characteristic
irreducibility claim, bounded actual inertia order, or finite
normalizer substitution has entered the argument.

\paragraph{Qualified source application.}
In the exact primitive minimal naturally nonsplit class of D-TD54,
the complete whole-source support floor is
\(M_0^{-1}\), where \(M_0=\max(80,N^*-2)\).
For a named unitary coarse layer, the accepted high-field
D-TD03 actual-factor argument excludes \(q\ge10^{12}\).
The theorem therefore gives a finite effective natural
dimension bound for every remaining unitary layer, including level
one, by taking the maximum of (UAR.3) over the finitely many prime
powers \(q<10^{12}\). For this application the whole floor is
applied directly to the original coefficient module and the same
complete D-TD12 extraction; an original subgroup lifting the later
projective core is not required. D-TD03 applies to the finite
special-linear lift of its actual projective factor image.
The sphere and \(r\ge12\) conditions follow from the actual
rank-nine source argument (NID.13), with \(e\ge2\) already forced
by (UAR.7). No extension to a more general original source class is
silently included. The separate D-TD69Q application records the
qualified quasiprimitive transfer with its own complete proof and review.
